% =====================================================================
% Section 5 -- Ex ante greening versus ex post insurance
% Paste into ehank_results.tex after the Experiments section (Sec. 4)
% and before Relation to the literature. All numbers below are the
% taste_shock = 0.05 / import-booking calibration; the \input-ed tables
% and figures regenerate them from run_summary_table.py,
% run_exante_expost.py and run_persistence.py.
% =====================================================================
\section{Ex ante greening versus ex post insurance}
\label{sec:exante}

The instruments in Section~\ref{sec:cap}--\ref{sec:transfer} are all
\emph{ex post}: they are deployed once the price shock has landed. Because the
adoption margin responds to the same relative price the crisis moves, a natural
alternative is to act \emph{ex ante}---to lean on the margin in the steady state
with a permanent carbon tax $\tau_b$ on brown energy, leaving the economy greener
and, one might hope, less exposed when the shock arrives. This section prices
that hope. The answer is that ex ante greening is a poor substitute for ex post
insurance, and the reason is intrinsic to the channel this paper adds: the
adoption margin is an endogenous shock absorber, and pre-greening spends it.

Table~\ref{tab:summary_price_import} collects the instruments on one axis. The
welfare metric is the total consumption-equivalent variation (CEV) of each full
scenario relative to the common no-tax steady state, by discount-factor type and
averaged, computed ex post from the cached impulse response
(\texttt{welfare.cev\_total}); a permanent policy enters through the standing gap
between its steady state and the baseline, a crisis-time policy through the
transition alone, so the two are directly comparable.

\input{output/tab_summary_import.tex}

\subsection{Targeting energy transfers is regressive}
\label{sec:flat}

The Slutsky transfer of Section~\ref{sec:transfer} indexes the lump sum to each
household's pre-crisis brown-energy use. Replacing it with an
\emph{untargeted} transfer of the \emph{same aggregate envelope}---the identical
outlay ($49.3$ vs.\ $49.3$ over $H=24$), paid flat across the distribution rather
than indexed---almost halves the welfare loss, from a CEV of $-0.82\%$ to
$-0.32\%$ (Table~\ref{tab:summary_price_import}). The gain is entirely
distributional. By type, the flat transfer moves the impatient (high-MPC)
household from $-0.68\%$ to $+0.10\%$---it is made \emph{better off} than in the
no-shock steady state---while the patient household bears more of the burden
($-0.77\%$); the Slutsky transfer, by contrast, spreads the loss almost evenly
($-0.68$, $-0.89$, $-0.88$).

The mechanism is that energy-indexed targeting is regressive \emph{in this
model}. With homothetic CES energy demand the energy share is constant across the
wealth distribution, so brown-energy use rises with total consumption, hence with
permanent income and patience. Indexing the transfer to pre-crisis energy use
therefore hands the largest cheques to the low-MPC, wealthy households, exactly
where a euro buys the least insurance. A flat transfer reaches the high-MPC
households the crisis hits hardest. This speaks directly to the design of the
2022--23 European energy-price ``brakes,'' which were indexed to historical
consumption: within the model, the indexing is the source of the inefficiency,
not a refinement of it.

We flag the caveat plainly. The result rests on the homothetic energy nest: if
energy were a necessity with a declining Engel curve---the empirically standard
case---low-income households would spend a \emph{larger} share on energy and the
sign could reverse. The clean statement is conditional: \emph{under constant
energy shares, energy-indexed targeting is dominated by an equal-value flat
transfer}. Making the nest non-homothetic is the right robustness check and is
left to the calibration section's agenda.

\subsection{A welfare-optimal permanent carbon tax}
\label{sec:carbontax}

Before comparing ex ante and ex post, we characterise the steady state a
permanent $\tau_b$ delivers. Holding $\psi_g$ fixed at its no-tax value and
letting $D^{G}_{ss}$ float, a carbon tax raises the effective brown/green price
ratio and lifts steady-state adoption: $D^{G}_{ss}$ moves from $5\%$ at
$\tau_b=0$ to $8.9\%$ at $\tau_b=0.10$ and $20.1\%$ at $\tau_b=0.30$
(Figure~\ref{fig:carbon_optimal}, right). The carbon revenue is rebated
lump-sum, so the exercise is budget-neutral and isolates the allocative effect.

Steady-state welfare is hump-shaped in $\tau_b$
(Figure~\ref{fig:carbon_optimal}, left): the standing CEV rises to a peak of
$+0.043\%$ at $\tau_b^{*}\approx 10\%$ and turns negative beyond
$\tau_b\approx 25\%$. The model contains \emph{no} climate externality, so this
interior optimum is not Pigouvian. It is a second-best correction of the
adoption friction: green energy is structurally cheaper
($\bar P^{E}_{G}=0.8<\bar P^{E}_{B}=1$) but the logit taste-shock scale
$\sigma_\varepsilon$ and the switching cost $\psi_g$ leave green under-adopted, so
a modest tax that pushes households toward the cheaper technology raises welfare
until the tax distortion dominates. That the optimum is positive at all is a
property of this friction, and we treat it as a diagnostic of the channel rather
than a policy recommendation---it would be swamped by any explicitly modelled
externality or abatement cost.

\begin{figure}[H]\centering
\includegraphics[width=\textwidth]{output/fig_carbon_optimal_import.pdf}
\caption{The welfare-optimal permanent carbon tax. Left: standing CEV of the
$\tau_b$ steady state relative to the no-tax baseline (no crisis), hump-shaped
with an interior optimum near $\tau_b^{*}\approx10\%$. Right: the induced
steady-state green share $D^{G}_{ss}$. $\psi_g$ is held fixed and $D^{G}_{ss}$
floats; the carbon revenue is rebated lump-sum.}
\label{fig:carbon_optimal}
\end{figure}

\subsection{Ex ante greening barely insures against the crisis}
\label{sec:versionA}

We now hit two economies with the \emph{same} price shock: the baseline
($D^{G}_{ss}=5\%$) and a ``prepared'' economy carrying the welfare-optimal
$\tau_b=0.10$ ($D^{G}_{ss}=8.9\%$). Table~\ref{tab:exante_expost_import}
reports the total CEV of each, alongside the ex post cap and transfer on the
baseline economy.

\input{output/tab_exante_expost_import.tex}

The prepared economy is essentially no better insured than the unprepared one:
its crisis CEV is $-1.85\%$ against $-1.87\%$ for laissez-faire, a difference of
two basis points, and both are far worse than either ex post instrument (cap
$-1.15\%$, transfer $-0.82\%$). Adding the standing benefit of the tax
($+0.04\%$, Section~\ref{sec:carbontax}) does not change the ranking. The
greening the tax buys is real---$D^{G}_{ss}$ is $78\%$ higher---but too small a
share of the economy to move aggregate exposure appreciably, and the tax's
crisis-time value is a rounding error next to a transfer that puts cash in
high-MPC hands exactly when it binds. Consistent with Section~\ref{sec:cap}, the
combination ``ETS $+$ cap'' tracks the cap alone: once the price is frozen the
adoption margin is shut and the greener starting point is irrelevant.

The green \emph{subsidy} of Section~\ref{sec:cap}\footnote{The transitory
switching subsidy $s_g$, financed by debt over the crisis window.} is worth a
line as a crisis instrument in its own right. It doubles the adoption response
(peak $D^{G}$ of $18.0$ vs.\ $8.9$ pp) at trivial fiscal cost, yet its CEV is
the worst in the table, $-2.03\%$, below laissez-faire. Paying households to
incur the real switching cost $\psi_g D^{sw}$---booked as an import
(Section~\ref{sec:bop})---in the middle of a consumption crisis destroys welfare
even as it greens the economy. The adoption margin is valuable, but not as
something to force at the trough.

\subsection{Persistence exhausts the adoption absorber}
\label{sec:routeA}

The obvious defence of the ex ante case is persistence: a longer-lived shock
should let the greener economy's lower exposure compound over more periods.
Figure~\ref{fig:persistence} and Table~\ref{tab:persistence_import} test it by
sweeping the shock half-life and re-solving only the impulse response (the steady
state does not move with persistence), decomposing the ex ante number into its
constant standing part and a \emph{crisis-only greening dividend}, defined as the
prepared economy's crisis CEV minus the standing $+0.04\%$ minus laissez-faire.

\input{output/tab_persistence_import.tex}

The defence fails, and in an instructive way. The greening dividend does not grow
with persistence---it \emph{turns negative}, from $+0.01\%$ at a four-quarter
half-life to $-0.09\%$ at a twelve-year one, crossing zero near a three-year
half-life. The ex post cap moves the other way: its crisis protection rises
monotonically, from $+0.24\%$ to $+1.36\%$, as a longer freeze shields more of
the transition.

The reason is the adoption margin itself. A persistent brown-price shock makes
switching attractive for many quarters, so the \emph{browner} laissez-faire
economy adopts heavily \emph{during} the crisis and reaps that windfall over the
whole episode---the margin is an endogenous shock absorber. The prepared economy
has already spent much of that absorber in the steady state: with fewer brown
households left at the switching margin, it has less adoption to do when the shock
arrives, and less of the crisis-time windfall to collect. Pre-greening does not
merely fail to insure against a persistent shock; it \emph{consumes the very
margin} that would have cushioned it. This is the sharpest form of the paper's
thesis. The instruments that freeze the adoption technology---the cap here, and
by construction the frozen-technology models of \citet{Bayer2026},
\citet{Langot2026} and \citet{ARS2024}---miss an absorber that works best when
left loaded, i.e.\ when the economy enters the crisis with room to green rather
than already green.

\begin{figure}[H]\centering
\includegraphics[width=\textwidth]{output/fig_persistence_import.pdf}
\caption{Persistence and the two policies. Left: total CEV of laissez-faire, the
ex post cap and the ex ante ETS as the shock half-life rises. Right: the
crisis-only gain over laissez-faire---the cap's protection grows with
persistence, while the ex ante \emph{greening dividend} turns negative, because a
more persistent shock lets the browner economy exploit the adoption absorber the
prepared economy has already spent.}
\label{fig:persistence}
\end{figure}

\paragraph{Takeaway.} Three results point the same way. Energy-indexed targeting
is dominated by a flat transfer of equal budget (Section~\ref{sec:flat}); a
permanent carbon tax is mildly welfare-improving in the steady state through a
non-Pigouvian adoption correction, but essentially worthless as crisis insurance
(Sections~\ref{sec:carbontax}--\ref{sec:versionA}); and its insurance value turns
negative as shocks become persistent, because greening ex ante exhausts the
adoption absorber (Section~\ref{sec:routeA}). The margin this paper adds is not a
lever a planner should pre-pull; it is a stabiliser that pays off precisely when
the economy is left free to adopt in the crisis.
